
\renewcommand{\arraystretch}{1.05}
\setlength{\parskip}{2pt}
% Extend only the figure into the margins: 19 cm on a 21 cm portrait A4 page.
\newcommand{\InlineUCModel}[1]{%
  \par\noindent
  \begin{minipage}{\linewidth}
    \centering
    \makebox[\linewidth][c]{%
      \includegraphics[width=19cm,height=0.60\textheight,keepaspectratio]{UC_#1.pdf}%
    }
    \captionsetup{hypcap=false}
    \captionof{figure}{UC-#1: detailed use-case model.}
    \label{fig:uc-model-#1}
  \end{minipage}
  \par
}
% ==================================================
% 2. USE-CASE SPECIFICATIONS
% ==================================================

\section{Use-case Models and Specifications}

This section describes the use cases of the
Medication Safety Assistant research prototype.
Each use case represents a user goal and an
observable result.

The use cases are organized by common user
workflows and supporting roles. Their order
does not imply a required execution sequence.
HCP reviews are simulated, and future
features are excluded.


% ==================================================
% USE-CASE OVERVIEW
% ==================================================

\subsection*{Use-Case Overview}

\begin{longtable}{|
    L{0.13\textwidth}|
    L{0.53\textwidth}|
    L{0.24\textwidth}|
}
\hline

\textbf{ID} &
\textbf{Use-case Name} &
\textbf{Actor(s)}
\\ \hline

\endfirsthead

\hline
\textbf{ID} &
\textbf{Use-case Name} &
\textbf{Actor(s)}
\\ \hline

\endhead

UC-01
& Sign In
& All registered roles
\\ \hline

UC-02
& Enter and Confirm Medication Information
& Patient
\\ \hline

UC-03
& Check Medication Safety
& Patient
\\ \hline

UC-04
& Request HCP Review
& Patient
\\ \hline

UC-05
& Review Medication Safety Case
& Healthcare Professional (HCP)
\\ \hline

UC-06
& Set Up Medication Reminders
& Patient
\\ \hline

UC-07
& Manage Medication Reminders
& Patient
\\ \hline

UC-08
& Record Dose Status
& Patient
\\ \hline

UC-09
& View Adherence History
& Patient
\\ \hline

UC-10
& View Medication and Safety History
& Patient
\\ \hline

UC-11
& Ask a Medication Question
& Patient
\\ \hline

UC-12
& Generate Patient Health Summary
& Patient
\\ \hline

UC-13
& Share Patient Health Summary with HCP
& Patient
\\ \hline

UC-14
& Review Medication Dataset Record
& Dataset Reviewer
\\ \hline

UC-15
& Manage User Accounts and Roles
& Administrator
\\ \hline

UC-16
& View Administrative Activity History
& Administrator
\\ \hline

UC-17
& Report Technical Issue
& All registered roles
\\ \hline

UC-18
& Handle Technical Support Request
& Technical Support
\\ \hline

\end{longtable}



% - 01. Access -

% UC-01 - Sign In | from UC-01
\clearpage
\subsection{Use-case: Sign In}
\label{uc:uc-01}
\InlineUCModel{01}

\begin{longtable}{|p{0.20\textwidth}|p{0.71\textwidth}|}
\hline
\textbf{Use case Name} & UC-01: Sign In \\ \hline
\textbf{Brief description} & A registered user signs in and receives access to the features permitted by their role. \\ \hline
\textbf{Actors} & Patient / General User; Healthcare Professional (HCP); Dataset Reviewer; Administrator; Technical Support \\ \hline
\textbf{Basic Flow} &
1. The user enters an account name and password.\par
2. The system checks the credentials and account status.\par
3. The system verifies the assigned role and permissions.\par
4. The system creates an authenticated session and opens the permitted dashboard.
\\ \hline
\textbf{Alternative Flows} &
\textbf{A1. Incorrect credentials (step 2)}\par
1. The system shows an error without creating a session.\par
2. The user may retry from step 1 or cancel.
\par
\textbf{A2. Inactive or blocked account (step 2)}\par
1. The system denies sign-in and advises the user to contact support.\par
2. The use case ends without a session.
\par
\textbf{A3. Role information unavailable (step 3)}\par
1. The system denies sign-in; no unrestricted session is created.\par
2. The use case ends with an access error.
\\* \hline
\textbf{Pre-conditions} & The user has a registered account and can open the sign-in page. \\* \hline
\textbf{Post-conditions} & \textbf{Success:} A session is created with the verified role permissions.\par
\textbf{Failure:} No authenticated session is created. \\ \hline
\end{longtable}

% - 02. Patient safety -

% UC-02 - Enter and Confirm Medication Information | from UC-02
\clearpage
\subsection{Use-case: Enter and Confirm Medication Information}
\label{uc:uc-02}
\InlineUCModel{02}

\begin{longtable}{|p{0.20\textwidth}|p{0.71\textwidth}|}
\hline
\textbf{Use case Name} & UC-02: Enter and Confirm Medication Information \\ \hline
\textbf{Brief description} & The patient enters medication information through a supported method and confirms the recorded details before further use. \\ \hline
\textbf{Actors} & Patient / General User \\ \hline
\textbf{Basic Flow} &
1. The patient provides medication information using OCR, chatbot text, voice, or manual input.\par
2. The system processes the input and identifies the requested task when needed.\par
3. The system shows the recorded medication details for review.\par
4. The patient corrects any errors and confirms that the details match the provided source.\par
5. The system saves the confirmed information, including any missing fields, for later use.
\\ \hline
\end{longtable}
\clearpage
\noindent\textbf{UC-02: Specification (continued)}
\begin{longtable}{|p{0.20\textwidth}|p{0.71\textwidth}|}
\hline
\textbf{Alternative Flows} &
\textbf{A1. Unclear request (step 2)}\par
1. The system asks for clarification or offers a function to select.\par
2. If clarified, return to step 2; otherwise, end the use case.
\par
\textbf{A2. Input processing fails (step 2)}\par
1. The system shows the error and offers retry or another input method.\par
2. If the patient retries, return to step 1; otherwise, end the use case.
\par
\textbf{A3. Missing or unclear details (step 3)}\par
1. The system marks the missing or unclear fields.\par
2. The patient corrects them when possible and continues to step 4; remaining gaps stay marked as incomplete.
\par
\textbf{A4. Patient does not confirm (step 4)}\par
1. The system does not release the input for safety checks or reminder activation.\par
2. The use case ends without confirmation.
\\* \hline
\textbf{Pre-conditions} & The patient is signed in and can access medication entry. \\* \hline
\textbf{Post-conditions} & \textbf{Success:} Confirmed input is saved; incomplete fields remain marked. Confirmation does not prove medication identity or safety.\par
\textbf{Failure:} Unconfirmed input is not used for safety checks or reminder activation. \\ \hline
\end{longtable}

% UC-03 - Check Medication Safety | from UC-03
\clearpage
\subsection{Use-case: Check Medication Safety}
\label{uc:uc-03}
\InlineUCModel{03}

\begin{longtable}{|p{0.20\textwidth}|p{0.71\textwidth}|}
\hline
\textbf{Use case Name} & UC-03: Check Medication Safety \\ \hline
\textbf{Brief description} & The patient receives supported drug-interaction and duplicate-ingredient findings, clear explanations, risk levels, data limits, and next steps without treatment changes. \\ \hline
\textbf{Actors} & Patient / General User \\ \hline
\textbf{Basic Flow} &
1. The patient requests a check using confirmed medication information.\par
2. The system attempts to identify the medicines and their active ingredients from supported data.\par
3. The system checks supported drug-drug interactions and duplicate active ingredients.\par
4. The system records risk levels for detected findings and the status of each supported or missing check.\par
5. The system shows affected medicines, explanations, warnings, data limitations, and appropriate next steps. When no issue is found, it limits that statement to completed checks.\par
6. The system saves the assessment, its status, and its time.
\\ \hline
\end{longtable}
\clearpage
\noindent\textbf{UC-03: Specification (continued)}
\begin{longtable}{|p{0.20\textwidth}|p{0.71\textwidth}|}
\hline
\textbf{Alternative Flows} &
\textbf{A1. Medicine cannot be identified (step 2)}\par
1. The system marks the affected medicine as unresolved and does not guess its ingredients.\par
2. It proceeds with assessable checks only; otherwise it reports an unassessed result at step 4.
\par
\textbf{A2. Required interaction data is missing (step 3)}\par
1. The system marks the affected interaction check unassessed and performs independent supported checks when possible.\par
2. Continue to step 4 with the limitation recorded.
\par
\textbf{A3. Severe interaction found (step 3)}\par
1. The system records and shows a severe warning immediately and recommends prompt HCP review.\par
2. The new-reminder safety gate holds activation only for affected medicines; existing reminders and prescriptions stay unchanged.\par
3. Continue to step 4 for remaining supported checks.
\par
\textbf{A4. Fewer than two identified medicines (step 3)}\par
1. If only one medicine was supplied, pairwise checks are not applicable.\par
2. If others are unresolved, their related checks remain unassessed; continue to step 4.
\par
\textbf{A5. Assessment fails (steps 2-4)}\par
1. The system reports a failed or incomplete assessment rather than a safe result.\par
2. Any severe warning already found and its reminder restriction remain recorded.\par
3. The patient may retry from step 1 or end the use case.
\\* \hline
\textbf{Pre-conditions} & The patient is signed in and has at least one confirmed medication record. \\* \hline
\textbf{Post-conditions} & \textbf{Success:} The assessment is saved as assessed, partially assessed, or unassessed, with findings and limitations. Severe interactions activate warnings and the new-reminder safety gate.\par
\textbf{Failure:} No failed or unassessed check is presented as safe; treatment and existing reminders remain unchanged. \\ \hline
\end{longtable}

% UC-04 - Request HCP Review | from UC-04
\clearpage
\subsection{Use-case: Request HCP Review}
\label{uc:uc-04}
\InlineUCModel{04}

\begin{longtable}{|p{0.20\textwidth}|p{0.71\textwidth}|}
\hline
\textbf{Use case Name} & UC-04: Request HCP Review \\ \hline
\textbf{Brief description} & The patient sends a selected medication safety case for simulated HCP review with permission to share the case. \\ \hline
\textbf{Actors} & Patient / General User \\ \hline
\textbf{Basic Flow} &
1. The patient chooses a saved safety case and selects Request HCP Review.\par
2. The system shows the case, findings, missing details, and information to be shared.\par
3. The patient reviews the information and gives permission to share the selected case.\par
4. The system creates a review request with Pending status and limits case access to authorized HCP accounts.\par
5. The system shows the request status to the patient.
\\ \hline
\end{longtable}
\clearpage
\noindent\textbf{UC-04: Specification (continued)}
\begin{longtable}{|p{0.20\textwidth}|p{0.71\textwidth}|}
\hline
\textbf{Alternative Flows} &
\textbf{A1. Patient declines sharing (step 3)}\par
1. The system creates no request or new sharing access.\par
2. The use case ends.
\par
\textbf{A2. Case is incomplete (step 2)}\par
1. The system highlights missing or unresolved details.\par
2. The patient may share the clearly marked limited case if minimum case information exists, or end the use case to update it.
\par
\textbf{A3. No reviewer is available (step 4)}\par
1. The request remains Pending without an assigned reviewer.\par
2. Continue to step 5; the system does not claim the case was reviewed.
\par
\textbf{A4. An active request already
exists (step 4)}\par
1. The system does not create
a duplicate request.\par
2. The system shows the existing
request and its current status.\par
3. No additional case information
is shared without permission.
The use case ends.
\par
\textbf{A5. Request cannot be saved (step 4)}\par
1. The system reports the failure and grants no new case access.\par
2. The patient may retry step 4 or end the use case.
\\* \hline
\textbf{Pre-conditions} & The patient is signed in and has a saved safety case. Sharing permission is requested within this use case. \\* \hline
\textbf{Post-conditions} & \textbf{Success:} A new Pending request, or an existing active request, is linked to the case. Only permitted HCP accounts can access the shared information.\par
\textbf{Failure:} No new request or sharing access is created; previous safety warnings and gates stay in force. \\ \hline
\end{longtable}

% - 03. HCP review -

% UC-05 - Review Medication Safety Case | from UC-14 + UC-15 + UC-16
\clearpage
\subsection{Use-case: Review Medication Safety Case}
\label{uc:uc-05}
\InlineUCModel{05}


\begin{longtable}{|p{0.20\textwidth}|p{0.71\textwidth}|}
\hline
\textbf{Use case Name} & UC-05: Review Medication Safety Case \\ \hline

\textbf{Brief description} &
An authorized HCP reviews a submitted safety case, communicates with the patient, and records a simulated outcome or requests missing information.
\\ \hline

\textbf{Actors} &
Healthcare Professional (HCP); Patient / General User
\\ \hline

\textbf{Basic Flow} &
1. The HCP opens the accessible review queue and selects a case awaiting a decision.\par
2. The system shows patient-approved case details, safety findings, warnings, data limits, and the existing conversation thread between the patient and HCP.\par
3. The HCP reviews the information and chooses a supported review outcome, with comments as needed.\par
4. The system checks the decision and records the reviewer, time, review outcome, and reply.\par
5. The system updates the case status and rechecks affected pending reminders when an eligible schedule confirmation is recorded.\par
6. The system shows the review outcome, HCP reply, and updated status to the patient and HCP.
\\ \hline

\end{longtable}
\clearpage
\noindent\textbf{UC-05: Specification (continued)}
\begin{longtable}{|p{0.20\textwidth}|p{0.71\textwidth}|}
\hline
\textbf{Alternative Flows} &

\textbf{A1. Case needs more information (step 3)}\par
1. The HCP specifies the missing details in the case conversation.\par
2. The system records the request and marks the case Awaiting Patient Response.\par
3. The patient submits the requested information in the same conversation thread.\par
4. The system marks the case Pending HCP Response and returns it to the HCP review queue.\par
5. When the HCP reopens the case, continue from step 2.

\par

\textbf{A2. Patient submits a follow-up question (after step 6)}\par
1. The patient submits a follow-up question in the existing case conversation.\par
2. The system appends the question to the same conversation thread and marks the case Pending HCP Response.\par
3. The system returns the case to the HCP review queue.\par
4. The HCP interface shows the chronological conversation as a reply thread and indicates that a new patient reply is available.\par
5. The HCP opens the case and continues from step 2.

\par

\textbf{A3. Case remains unresolved (step 3)}\par
1. The HCP selects Unresolved and records the reason.\par
2. Continue to step 4 to save the review outcome. Required reminder holds remain in place.

\par

\textbf{A4. No valid schedule confirmation (step 5)}\par
1. Affected new reminders remain Pending; no prescribed timing is changed.\par
2. Continue to step 6 with the remaining restriction displayed.

\par

\textbf{A5. Case is unauthorized, changed, or save fails (steps 1 or 4)}\par
1. The system denies or stops the action and reports the issue.\par
2. No unrecorded decision or reply is treated as approval.
\\* \hline
\textbf{Pre-conditions} &
The HCP is signed in with access to the selected request; the case was shared with patient permission.
\\* \hline
\textbf{Post-conditions} &
\textbf{Success:} A simulated, traceable review outcome, information request, or conversation reply is recorded. A patient follow-up returns the case to the permitted HCP queue. Only an eligible, recorded confirmation can affect pending new reminders.\par
\textbf{Failure:} No unrecorded decision changes access, prescriptions, or reminder restrictions.
\\ \hline

\end{longtable}

% - 04. Reminders and adherence -

% UC-06 - Set Up Medication Reminders | from UC-07
\clearpage
\subsection{Use-case: Set Up Medication Reminders}
\label{uc:uc-06}
\InlineUCModel{06}

\begin{longtable}{|p{0.20\textwidth}|p{0.71\textwidth}|}
\hline
\textbf{Use case Name} & UC-06: Set Up Medication Reminders \\ \hline
\textbf{Brief description} & The patient activates reminders from an existing, confirmed medication schedule when the applicable safety rules allow it. \\ \hline
\textbf{Actors} & Patient / General User \\ \hline
\textbf{Basic Flow} &
1. The patient selects a medication with existing prescription or recorded HCP instructions.\par
2. The system displays the recorded schedule, its source, and relevant safety status.\par
3. The patient confirms that the schedule matches the existing instructions.\par
4. The system checks schedule verification, recorded risk warnings, and the new-reminder safety gate.\par
5. The system activates eligible reminders and shows their status.
\\ \hline
\end{longtable}
\clearpage
\noindent\textbf{UC-06: Specification (continued)}
\begin{longtable}{|p{0.20\textwidth}|p{0.71\textwidth}|}
\hline
\textbf{Alternative Flows} &
\textbf{A1. Schedule is missing or unverified (steps 2-3)}\par
1. The system does not activate a reminder without a verified schedule.\par
2. The patient may correct the information or end the use case.
\par
\textbf{A2. Patient proposes new timing (step 3)}\par
1. The system records a proposed change as Pending HCP Review.\par
2. It does not activate that proposed schedule before the required simulated confirmation.
\par
\textbf{A3. Severe interaction affects the medicine (step 4)}\par
1. The system keeps affected new reminders Pending HCP Review.\par
2. Other eligible reminders may activate; existing reminders are not stopped automatically.
\par
\textbf{A4. Safety check is unavailable (step 4)}\par
1. The system reports Unassessed rather than safe clearance.\par
2. The reminder remains restricted when the applicable review policy requires confirmation.
\par
\textbf{A5. Browser notifications are unavailable (step 5)}\par
1. The system reports the notification limitation and still shows eligible schedules in the application.
\\* \hline
\textbf{Pre-conditions} & The patient is signed in, and a medication record contains instructions that can be verified. \\* \hline
\textbf{Post-conditions} & \textbf{Success:} Eligible new reminders are activated; restricted ones remain Pending with the reason recorded.\par
\textbf{Failure:} No unverified or blocked new reminder is activated; prescriptions and existing reminders remain unchanged. \\ \hline
\end{longtable}

% UC-07 - Manage Medication Reminders | from UC-08
\clearpage
\subsection{Use-case: Manage Medication Reminders}
\label{uc:uc-07}
\InlineUCModel{07}

\begin{longtable}{|p{0.20\textwidth}|p{0.71\textwidth}|}
\hline
\textbf{Use case Name} & UC-07: Manage Medication Reminders \\ \hline
\textbf{Brief description} & The patient changes permitted reminder settings or disables notifications without changing prescribed treatment. \\ \hline
\textbf{Actors} & Patient / General User \\ \hline
\textbf{Basic Flow} &
1. The patient opens an existing reminder and chooses to manage it.\par
2. The system shows the current schedule, its source, and notification settings.\par
3. The patient requests a notification change or a schedule update.\par
4. The system checks whether the change is allowed by the verified instructions and safety rules.\par
5. The patient confirms the request.\par
6. The system saves an allowed setting or records a proposed timing change as Pending HCP Review.
\\ \hline
\textbf{Alternative Flows} &
\textbf{A1. Patient disables notifications (step 3)}\par
1. The system prepares to disable
only the selected notification.\par
2. Continue to step 5 for confirmation.
The medication instructions remain unchanged.
\par

\textbf{A2. Timing differs from verified
instructions (step 4)}\par
1. The system marks the proposed
change as requiring HCP review.\par
2. Continue to step 5 for patient
confirmation. The existing schedule
remains unchanged.

\par
\textbf{A3. Patient cancels or update fails (steps 5-6)}\par
1. The system keeps the previous reminder settings and reports an error when relevant.
\\* \hline
\textbf{Pre-conditions} & The patient is signed in and owns an existing reminder. \\* \hline
\textbf{Post-conditions} & \textbf{Success:} A permitted reminder setting is saved, or a timing proposal is Pending review.\par
\textbf{Failure:} The prior reminder settings remain effective unless the patient explicitly disabled a notification. \\ \hline
\end{longtable}

% UC-08 - Record Dose Status | from UC-09
\clearpage
\subsection{Use-case: Record Dose Status}
\label{uc:uc-08}
\InlineUCModel{08}

\begin{longtable}{|p{0.20\textwidth}|p{0.71\textwidth}|}
\hline
\textbf{Use case Name} & UC-08: Record Dose Status \\ \hline
\textbf{Brief description} & The patient records whether a scheduled medication dose was Taken, Missed, or Skipped for adherence tracking. \\ \hline
\textbf{Actors} & Patient / General User \\ \hline
\textbf{Basic Flow} &
1. The patient selects a scheduled dose from the reminder or adherence view.\par
2. The system shows the medicine, scheduled time, and current dose status.\par
3. The patient selects Taken, Missed, or Skipped.\par
4. The patient confirms the entry.\par
5. The system saves the dose status and recorded time.
\\ \hline
\textbf{Alternative Flows} &
\textbf{A1. Existing entry is corrected (step 3)}\par
1. The system shows the previous status for comparison.\par
2. The patient confirms the correction; continue to step 5.
\par
\textbf{A2. Patient cancels (step 4)}\par
1. The system keeps the previous entry and ends the use case.
\par
\textbf{A3. Dose event is unavailable or save fails (steps 2 or 5)}\par
1. The system explains that the entry cannot be saved.\par
2. No unsupported dose record is created.
\\* \hline
\textbf{Pre-conditions} & The patient is signed in and has an accessible scheduled dose or adherence entry. \\* \hline
\textbf{Post-conditions} & \textbf{Success:} The confirmed dose status is saved and available in adherence history.\par
\textbf{Failure:} No dose entry is changed; no treatment recommendation is produced. \\ \hline
\end{longtable}

% UC-09 - View Adherence History | from UC-10
\clearpage
\subsection{Use-case: View Adherence History}
\label{uc:uc-09}
\InlineUCModel{09}

\begin{longtable}{|p{0.20\textwidth}|p{0.71\textwidth}|}
\hline
\textbf{Use case Name} & UC-09: View Adherence History \\ \hline
\textbf{Brief description} & The patient reviews recorded medication use over time. \\ \hline
\textbf{Actors} & Patient / General User \\ \hline
\textbf{Basic Flow} &
1. The patient opens Adherence History.\par
2. The system shows available scheduled doses and recorded statuses.\par
3. The patient selects a medicine or time period.\par
4. The system displays matching Taken, Missed, and Skipped records with dates.
\\ \hline
\textbf{Alternative Flows} &
\textbf{A1. No entries match (step 4)}\par
1. The system shows an empty result and allows another filter.
\par
\textbf{A2. History is incomplete or unavailable (step 2)}\par
1. The system displays available records and reports any unavailable period.
\\* \hline
\textbf{Pre-conditions} & The patient is signed in and can view their own adherence records. \\* \hline
\textbf{Post-conditions} & \textbf{Success:} Permitted adherence history is displayed without modifying it.\par
\textbf{Failure:} No record is changed. \\ \hline
\end{longtable}

% - 05. Records and history -

% UC-10 - View Medication and Safety History | from UC-06
\clearpage
\subsection{Use-case: View Medication and Safety History}
\label{uc:uc-10}
\InlineUCModel{10}

\begin{longtable}{|p{0.20\textwidth}|p{0.71\textwidth}|}
\hline
\textbf{Use case Name} & UC-10: View Medication and Safety History \\ \hline
\textbf{Brief description} & The patient reviews saved medication records, past safety results, review outcomes, and their dates. \\ \hline
\textbf{Actors} & Patient / General User \\ \hline
\textbf{Basic Flow} &
1. The patient opens Medication and Safety History.\par
2. The system lists permitted medication records and past safety assessments.\par
3. The patient selects a record or assessment.\par
4. The system shows its date, findings, limitations, and linked HCP review status.\par
5. The system clearly labels results that no longer match the current medication record.
\\ \hline
\textbf{Alternative Flows} &
\textbf{A1. No history is available (step 2)}\par
1. The system displays an empty-state message and ends the use case.
\par
\textbf{A2. A record is no longer available or access is denied (step 3)}\par
1. The system reports the limitation without showing private data.\par
2. The patient may select another permitted record.
\\* \hline
\textbf{Pre-conditions} & The patient is signed in and can view their own records. \\* \hline
\textbf{Post-conditions} & \textbf{Success:} Historical information is displayed without changing stored data.\par
\textbf{Failure:} No record or treatment is changed. \\ \hline
\end{longtable}

% - 06. Chatbot and summaries -

% UC-11 
\clearpage
\subsection{Use-case: Ask a Medication Question}
\label{uc:uc-11}
\InlineUCModel{11}

\begin{longtable}{|p{0.20\textwidth}|p{0.71\textwidth}|}
\hline
\textbf{Use case Name} & UC-11: Ask a Medication Question \\ \hline
\textbf{Brief description} & The patient receives a supported medication-related answer from the chatbot with appropriate limits. \\ \hline
\textbf{Actors} & Patient / General User \\ \hline
\textbf{Basic Flow} &
1. The patient asks a medication-related question in text or voice.\par
2. The system identifies the request and any relevant authorized context.\par
3. The system checks supported reference information and permitted patient records.\par
4. The system answers with relevant explanations and limits or asks a needed follow-up.\par
5. The patient may ask another question or open a related supported feature.
\\ \hline
\textbf{Alternative Flows} &
\textbf{A1. Question is unclear (step 2)}\par
1. The system asks a short follow-up question.\par
2. If clarified, return to step 2; otherwise, end the use case.
\par
\textbf{A2. Reliable information is unavailable (step 3)}\par
1. The system explains the information limit rather than inventing medical claims.\par
2. The use case ends or the patient asks another question.
\par
\textbf{A3. Request is outside scope or suggests urgent danger (step 2)}\par
1. The chatbot does not diagnose, prescribe, change treatment, or provide unvalidated emergency instructions.\par
2. It gives an appropriate warning to seek qualified or emergency help.
\par
\textbf{A4. Request would modify records (step 5)}\par
1. The system opens the relevant separate confirmation use case.\par
2. No data is changed solely from an unconfirmed chatbot message.
\\* \hline
\textbf{Pre-conditions} & The patient is signed in and the chatbot is available. \\* \hline
\textbf{Post-conditions} & \textbf{Success:} The patient receives a supported answer, clarification, or clear limitation.\par
\textbf{Failure:} No unsupported clinical decision or unconfirmed record change occurs. \\ \hline
\end{longtable}

% UC-12 - Generate Patient Health Summary | from UC-12
\clearpage
\subsection{Use-case: Generate Patient Health Summary}
\label{uc:uc-12}
\InlineUCModel{12}

\begin{longtable}{|p{0.20\textwidth}|p{0.71\textwidth}|}
\hline
\textbf{Use case Name} & UC-12: Generate Patient Health Summary \\ \hline
\textbf{Brief description} & The patient creates a structured, reviewable summary from permitted medication and reported health information. \\ \hline
\textbf{Actors} & Patient / General User \\ \hline
\textbf{Basic Flow} &
1. The patient requests a Patient Health Summary.\par
2. The system gathers permitted medication records, patient-reported details, and available safety assessments.\par
3. The system generates a draft summary with dates and known information gaps.\par
4. The patient reviews the draft and corrects inaccurate reported details.\par
5. The patient confirms the summary; the system saves it without sharing it automatically.
\\ \hline
\textbf{Alternative Flows} &
\textbf{A1. Information is missing (step 2)}\par
1. The system labels unavailable items and prepares only a limited draft.\par
2. Continue to step 3.
\par
\textbf{A2. The draft needs correction (step 4)}\par
1. The patient edits supported fields or corrects source records.\par
2. The system updates the draft and returns to step 4.
\par
\textbf{A3. Patient cancels or saving fails (step 5)}\par
1. No new confirmed summary is saved or shared.\par
2. The patient may retry if saving failed.
\\* \hline
\textbf{Pre-conditions} & The patient is signed in and has at least one permitted record or reported health detail. \\* \hline
\textbf{Post-conditions} & \textbf{Success:} A patient-reviewed summary is stored for later viewing or optional sharing; it is not a diagnosis.\par
\textbf{Failure:} No unconfirmed summary is shared. \\ \hline
\end{longtable}

% UC-13 - Share Patient Health Summary with HCP | from UC-13
\clearpage
\subsection{Use-case: Share Patient Health Summary with HCP}
\label{uc:uc-13}
\InlineUCModel{13}

\begin{longtable}{|p{0.20\textwidth}|p{0.71\textwidth}|}
\hline
\textbf{Use case Name} & UC-13: Share Patient Health Summary with HCP \\ \hline
\textbf{Brief description} & The patient grants an authorized HCP access to a selected, confirmed health summary. \\ \hline
\textbf{Actors} & Patient / General User \\ \hline
\textbf{Basic Flow} &
1. The patient selects a saved summary and chooses Share with HCP.\par
2. The system shows the summary and the intended authorized recipient or review case.\par
3. The patient reviews the scope and confirms permission to share.\par
4. The system grants access only to the selected summary for that recipient or case.\par
5. The system shows the sharing status.
\\ \hline
\textbf{Alternative Flows} &
\textbf{A1. Patient declines permission (step 3)}\par
1. No new access is granted; the use case ends.
\par
\textbf{A2. Recipient or case is unauthorized (step 4)}\par
1. The system denies sharing and asks for an authorized choice.\par
2. The patient may return to step 2 or cancel.
\par
\textbf{A3. Summary is outdated (step 2)}\par
1. The system warns the patient about its date or missing information.\par
2. The patient may share it with its limits visible or return to Generate Patient Health Summary.
\par
\textbf{A4. Sharing fails (step 4)}\par
1. The system reports the error and does not claim success or create unintended access.
\\* \hline
\textbf{Pre-conditions} & The patient is signed in and has a saved, reviewed summary; an authorized target can be selected. \\* \hline
\textbf{Post-conditions} & \textbf{Success:} Only the chosen summary is accessible to its permitted HCP recipient or case.\par
\textbf{Failure:} No new access is granted. \\ \hline
\end{longtable}

% - 07. Dataset review -

% UC-14 - Review Medication Dataset Record | from UC-17 + UC-18
\clearpage
\subsection{Use-case: Review Medication Dataset Record}
\label{uc:uc-14}
\InlineUCModel{14}

\begin{longtable}{|p{0.20\textwidth}|p{0.71\textwidth}|}
\hline
\textbf{Use case Name} & UC-14: Review Medication Dataset Record \\ \hline
\textbf{Brief description} & The Dataset Reviewer checks a medication record against its evidence and records a traceable quality decision. \\ \hline
\textbf{Actors} & Dataset Reviewer \\ \hline
\textbf{Basic Flow} &
1. The reviewer opens the permitted dataset review list and selects a medication record.\par
2. The system shows names, ingredients, mappings, source references, and review status.\par
3. The reviewer compares the record with the available evidence.\par
4. The reviewer selects Accept, Reject, or Needs Correction and enters a reason.\par
5. The system checks the required review details and saves the decision, reviewer, time, and evidence.\par
6. The system updates the dataset review status.
\\ \hline
\end{longtable}
\clearpage
\noindent\textbf{UC-14: Specification (continued)}
\begin{longtable}{|p{0.20\textwidth}|p{0.71\textwidth}|}
\hline
\textbf{Alternative Flows} &
\textbf{A1. Source evidence is missing (steps 2-4)}\par
1. The reviewer selects Needs Correction or Reject and states the missing evidence.\par
2. The system does not mark the record as clinically validated.
\par
\textbf{A2. Inconsistent mapping is found (step 3)}\par
1. The reviewer records the issue and chooses Reject or Needs Correction.\par
2. Continue to step 5.
\par
\textbf{A3. Record changed during review (step 5)}\par
1. The system asks the reviewer to reload the latest version before deciding.\par
2. Return to step 2.
\par
\textbf{A4. Reviewer cancels or saving fails (steps 4-5)}\par
1. No new decision is saved; the previous review status remains.
\\* \hline
\textbf{Pre-conditions} & The Dataset Reviewer is signed in and permitted to review the record. \\* \hline
\textbf{Post-conditions} & \textbf{Success:} A source-linked data-quality decision and change history are recorded; no independent clinical validation is claimed.\par
\textbf{Failure:} No new review decision is saved. \\ \hline
\end{longtable}

% - 08. Administration -

% UC-15 - Manage User Accounts and Roles | from UC-19 + UC-20
\clearpage
\subsection{Use-case: Manage User Accounts and Roles}
\label{uc:uc-15}
\InlineUCModel{15}

\begin{longtable}{|p{0.20\textwidth}|p{0.71\textwidth}|}
\hline
\textbf{Use case Name} & UC-15: Manage User Accounts and Roles \\ \hline
\textbf{Brief description} & The Administrator creates or updates accounts and assigns permitted roles under the system access policy. \\ \hline
\textbf{Actors} & Administrator \\ \hline
\textbf{Basic Flow} &
1. The Administrator opens account management and selects an existing account or Create Account.\par
2. The system displays the account fields, status, and permitted role options.\par
3. The Administrator enters account details, a status change, or a role assignment.\par
4. The system validates the request against account and RBAC rules.\par
5. The Administrator confirms the proposed change.\par
6. The system saves the authorized change and records the administrative activity.
\\ \hline
\textbf{Alternative Flows} &
\textbf{A1. Required data is invalid or account already exists (step 4)}\par
1. The system explains the problem and requests correction.\par
2. Return to step 3, or end if canceled.
\par
\textbf{A2. Role change or action is not permitted (step 4)}\par
1. The system denies the request without granting extra permissions.\par
2. The use case ends or the Administrator selects another valid action.
\par
\textbf{A3. Administrator cancels or save fails (steps 5-6)}\par
1. The previous account and role settings remain unchanged.\par
2. The system reports a save error when relevant.
\\* \hline
\textbf{Pre-conditions} & The Administrator is signed in with account and role-management permission. \\* \hline
\textbf{Post-conditions} & \textbf{Success:} The permitted account or role change is stored with an activity record.\par
\textbf{Failure:} No unauthorized change is saved. \\ \hline
\end{longtable}

% UC-16 - View Administrative Activity History | from UC-21
\clearpage
\subsection{Use-case: View Administrative Activity History}
\label{uc:uc-16}
\InlineUCModel{16}

\begin{longtable}{|p{0.20\textwidth}|p{0.71\textwidth}|}
\hline
\textbf{Use case Name} & UC-16: View Administrative Activity History \\ \hline
\textbf{Brief description} & The Administrator examines permitted audit records for account and access-management actions. \\ \hline
\textbf{Actors} & Administrator \\ \hline
\textbf{Basic Flow} &
1. The Administrator opens Administrative Activity History.\par
2. The system shows permitted activity records and their dates.\par
3. The Administrator filters by date, account, or action type.\par
4. The system shows matching events and responsible accounts without restricted patient content.
\\ \hline
\textbf{Alternative Flows} &
\textbf{A1. No records match (step 4)}\par
1. The system shows an empty result and allows a different filter.
\par
\textbf{A2. History service is unavailable (step 2)}\par
1. The system reports the error and does not invent activity records.
\\* \hline
\textbf{Pre-conditions} & The Administrator is signed in and permitted to view audit metadata. \\* \hline
\textbf{Post-conditions} & \textbf{Success:} Permitted activity history is displayed without changing the records.\par
\textbf{Failure:} No audit or patient data is changed. \\ \hline
\end{longtable}

% - 09. Technical support -

% UC-17 - Report Technical Issue | from UC-22
\clearpage
\subsection{Use-case: Report Technical Issue}
\label{uc:uc-17}
\InlineUCModel{17}

\begin{longtable}{|p{0.20\textwidth}|p{0.71\textwidth}|}
\hline
\textbf{Use case Name} & UC-17: Report Technical Issue \\ \hline
\textbf{Brief description} & A registered user reports a technical problem and receives a support ticket status. \\ \hline
\textbf{Actors} & Patient / General User; Healthcare Professional (HCP); Dataset Reviewer; Administrator; Technical Support \\ \hline
\textbf{Basic Flow} &
1. The user opens Report Technical Issue.\par
2. The system shows the required issue fields and optional diagnostic details.\par
3. The user describes the problem and reviews any attached information.\par
4. The user submits the report.\par
5. The system saves a ticket and shows its identifier and status.
\\ \hline
\textbf{Alternative Flows} &
\textbf{A1. Required information is missing (step 4)}\par
1. The system highlights the missing information.\par
2. Return to step 3 for correction.
\par
\textbf{A2. Attachment contains unnecessary patient information (step 3)}\par
1. The system requests removal or masking of the private data before submission.\par
2. Return to step 3.
\par
\textbf{A3. User cancels or ticket creation fails (steps 4-5)}\par
1. No successful ticket submission is reported.\par
2. The user may retry after an error or end the use case.
\\* \hline
\textbf{Pre-conditions} & The user is signed in and can open the technical support feature. \\* \hline
\textbf{Post-conditions} & \textbf{Success:} A support ticket is stored and its identifier is shown without unnecessary patient data.\par
\textbf{Failure:} No ticket is created. \\ \hline
\end{longtable}

% UC-18 - Handle Technical Support Request | from UC-23
\clearpage
\subsection{Use-case: Handle Technical Support Request}
\label{uc:uc-18}
\InlineUCModel{18}

\begin{longtable}{|p{0.20\textwidth}|p{0.71\textwidth}|}
\hline
\textbf{Use case Name} & UC-18: Handle Technical Support Request \\ \hline
\textbf{Brief description} & Technical Support investigates a permitted ticket, records progress, and provides an outcome when available. \\ \hline
\textbf{Actors} & Technical Support \\ \hline
\textbf{Basic Flow} &
1. Support staff opens the permitted ticket queue and selects a ticket.\par
2. The system shows the issue description, status, and allowed technical details.\par
3. Support staff investigates the issue and records relevant notes.\par
4. Support staff records a resolution or updates the ticket status.\par
5. The system saves the action, time, and responsible account and shows the updated status to the reporting user.
\\ \hline
\end{longtable}
\clearpage
\noindent\textbf{UC-18: Specification (continued)}
\begin{longtable}{|p{0.20\textwidth}|p{0.71\textwidth}|}
\hline
\textbf{Alternative Flows} &
\textbf{A1. More information is required (step 3)}\par
1. Support staff requests only necessary technical details.\par
2. The ticket becomes Pending User Response; the current use case ends.
\par
\textbf{A2. The issue is not resolved (step 4)}\par
1. Support staff records progress and keeps the ticket Open or In Progress.\par
2. Continue to step 5.
\par
\textbf{A3. Unauthorized patient data is present (step 2)}\par
1. The system restricts access to medical content.\par
2. Support staff works with permitted technical details only.
\par
\textbf{A4. Access is denied or saving fails (steps 1 or 5)}\par
1. The system denies the action or reports the error.\par
2. No unsupported resolution is recorded.
\\* \hline
\textbf{Pre-conditions} & Support staff is signed in and permitted to access the selected support ticket. \\* \hline
\textbf{Post-conditions} & \textbf{Success:} The ticket status and troubleshooting history are saved; the reporter can view the status.\par
\textbf{Failure:} No unsupported ticket update or patient-data access occurs. \\ \hline
\end{longtable}
